% ============================================================
% cours_fdtd_eme_fem.tex — DOCUMENT MAÎTRE
%
% Cours : Simulation électromagnétique (FDTD/EME/FEM) et
% interaction avec les fonderies photoniques.
%
% Ce fichier est LE SEUL point d'entrée de compilation. Il porte
% le préambule, la page de garde et la table des matières
% uniques du cours. Chaque module est un fichier séparé sous
% modules/ inséré via \include, afin que la compilation ne
% produise jamais qu'un unique PDF cohérent.
%
% Compilation (XeLaTeX, deux passes minimum, trois si la table
% des matières ou les références changent) :
%   xelatex cours_fdtd_eme_fem.tex
%   xelatex cours_fdtd_eme_fem.tex
% ============================================================

\documentclass[11pt,twoside,openright]{book}

\usepackage{style_cours}

% ---------- Métadonnées PDF ----------
\hypersetup{
  pdftitle={Simulation photonique : FDTD, EME, FEM et interaction fonderie},
  pdfauthor={},
  pdfsubject={Cours de simulation électromagnétique en photonique intégrée},
  colorlinks=true,
  linkcolor=coursbleu,
  citecolor=coursbleu,
  urlcolor=coursbleuclair,
  bookmarksnumbered=true,
  bookmarksdepth=3,
}

% ---------- Compteur de statut d'unité (usage interne) ----------
% \'Etat de rédaction affiché en filigrane de chapitre tant qu'une
% unité n'a que son plan (PLAN) ou son contenu complet (COMPLET).
% Purement informatif ; à retirer manuellement quand toutes les
% unités du cours sont passées en COMPLET.
\newcommand{\etatunite}[1]{%
  \begin{flushright}
    \small\itshape\textcolor{coursgris}{[État de rédaction : #1]}
  \end{flushright}
}

\begin{document}

% ============================================================
% PAGE DE GARDE (unique pour tout le cours)
% ============================================================
\begin{titlepage}
  \centering
  \vspace*{2cm}

  {\color{coursbleu}\rule{\textwidth}{1.5pt}}\\[0.4cm]
  {\Huge\bfseries\color{coursbleu} Simulation électromagnétique intégrée\par}
  \vspace{0.3cm}
  {\Large\bfseries\color{coursbleuclair} FDTD, EME, FEM et interaction avec les fonderies photoniques\par}
  \vspace{0.4cm}
  {\color{coursbleu}\rule{\textwidth}{1.5pt}}\\[1.2cm]

  {\large\itshape De l'électromagnétisme guidé au tape-out~:\\
   un parcours code-first pour physicien expérimenté\par}

  \vspace{2.5cm}
  \begin{tcolorbox}[base,colback=coursbleu!5,colframe=coursbleu,width=0.75\textwidth]
    \textbf{Public visé~:} physicien(ne) titulaire d'un doctorat en optique
    guidée / physique laser, disposant d'une base théorique solide en
    électromagnétisme mais découvrant les outils de simulation industriels
    (Lumerical, IPKISS, RSoft) et le vocabulaire fonderie.
  \end{tcolorbox}

  \vfill

  {\large\today\par}
  \vspace{0.5cm}
  {\normalsize Document maître unique — compilé en XeLaTeX\par}
\end{titlepage}

\frontmatter

\tableofcontents

\mainmatter

% ============================================================
% MODULE 0 — Rappels fondamentaux
% ============================================================
\part{Rappels fondamentaux~: optique guidée, électromagnétisme appliqué, plateformes matériaux}
\label{part:module0}
% ============================================================
% Module 0 — Rappels fondamentaux
% ÉTAT : PLAN (titres, objectifs, prérequis, outils, aperçu).
% Le contenu pédagogique complet (théorie développée, rappels
% rédigés, simulations guidées, exercices) sera généré chapitre
% par chapitre dans une session ultérieure, après validation de
% ce plan.
% ============================================================

\chapter{Optique guidée et modes~: rappel opérationnel}
\label{chap:m0-optique}
\etatunite{PLAN}

\noindent\textbf{Niveau~:} débutant (dans le contexte industriel) \hfill
\textbf{Position~:} Module 0, Chapitre 1 sur 4

\begin{objectifs}
\begin{itemize}[leftmargin=1.4em]
  \item \textbf{[Théorique]} Reformuler la notion d'indice effectif $n_{\text{eff}}$ et de constante de propagation $\beta$ à partir de l'équation de dispersion d'un guide plan.
  \item \textbf{[Théorique]} Distinguer modes guidés, modes de fuite (leaky) et modes de substrat/radiation, et relier ce vocabulaire à celui, déjà connu, des modes de cavité laser.
  \item \textbf{[Théorique]} Caractériser la biréfringence de forme (TE/TM) propre aux guides rectangulaires intégrés, absente des fibres à symétrie cylindrique.
  \item \textbf{[Pratique]} Calculer numériquement $n_{\text{eff}}$ d'un guide plan à 3 couches en Python (méthode graphique/résolution de l'équation transcendante), sans solveur dédié.
  \item \textbf{[Pratique]} Identifier, sur un profil de mode tracé, les grandeurs qui seront reprises tout au long du cours (largeur à mi-hauteur, confinement, aire effective).
\end{itemize}
\end{objectifs}

\begin{prerequis}
\textbf{Théoriques~:} équations de Maxwell en milieu diélectrique linéaire, notion de mode propre, réflexion totale interne — acquis de thèse, aucun renvoi interne nécessaire.\\
\textbf{Outils/logiciels~:} Python 3, NumPy, Matplotlib (environnement à installer si non déjà disponible).
\end{prerequis}

\noindent\textbf{Outils principaux~:} Python/NumPy (open source, aucune dépendance photonique spécifique à ce stade).

\noindent\textbf{Rappel de mise à niveau prévu~:} non — ce chapitre \emph{est} lui-même le rappel d'ancrage du cours ; il n'y a pas de notion industrielle nouvelle encore introduite.

\noindent\textbf{Aperçu du contenu~:} pont explicite entre le formalisme \og modes de cavité / faisceaux gaussiens\fg{} déjà maîtrisé et le formalisme \og modes de guide planaire/rectangulaire\fg{} utilisé par tous les solveurs du cours (FDTD/EME/FEM). Introduction du vocabulaire anglais standard (\emph{effective index}, \emph{leaky mode}, \emph{form birefringence}) qui réapparaîtra sans retraduction dans les modules suivants.

\noindent\textbf{Livrable prévu~:} script Python autonome traçant $n_{\text{eff}}(\lambda)$ pour un guide plan symétrique et asymétrique, réutilisé comme cas de validation dans le Module 1.

\bigskip
\hrule
\bigskip

\chapter{Électromagnétisme appliqué~: des équations de Maxwell aux formulations numériques}
\label{chap:m0-em}
\etatunite{PLAN}

\noindent\textbf{Niveau~:} débutant (dans le contexte industriel) \hfill
\textbf{Position~:} Module 0, Chapitre 2 sur 4

\begin{objectifs}
\begin{itemize}[leftmargin=1.4em]
  \item \textbf{[Théorique]} Réécrire les équations de Maxwell sous les trois formes qui seront respectivement à la base de FDTD (formulation temporelle, rotationnels couplés E/H), EME (formulation modale/harmonique, propagation par matrices S) et FEM (formulation faible/variationnelle).
  \item \textbf{[Théorique]} Expliquer pourquoi un même problème physique (mode d'un guide) admet trois formulations numériques différentes et quels compromis (temps, mémoire, généralité géométrique) en découlent.
  \item \textbf{[Théorique]} Relier la notion de condition aux limites absorbante (PML) à un concept déjà connu~: l'apodisation/l'absorption en bord de cavité pour éviter les réflexions parasites.
  \item \textbf{[Pratique]} Identifier dans un fichier de configuration Meep/Lumerical à quelle équation/formulation correspond chaque bloc (source, milieu, moniteur, condition aux limites).
\end{itemize}
\end{objectifs}

\begin{prerequis}
\textbf{Théoriques~:} \cref{chap:m0-optique} (modes, $n_{\text{eff}}$); équations de Maxwell macroscopiques — acquis de thèse.\\
\textbf{Outils/logiciels~:} aucun nouveau ; lecture de fichiers de configuration à titre d'illustration uniquement (pas encore d'exécution).
\end{prerequis}

\noindent\textbf{Outils principaux~:} aucun solveur exécuté dans ce chapitre — chapitre de mise en carte conceptuelle avant l'entrée en matière outillée du Module~1.

\noindent\textbf{Rappel de mise à niveau prévu~:} oui, court encart sur les conditions aux limites absorbantes (PML), formulé comme extension du concept d'apodisation de cavité déjà connu, avant sa réapparition opérationnelle en \cref{chap:m0-optique} et au Module~1.

\noindent\textbf{Aperçu du contenu~:} tableau de correspondance \{formulation mathématique $\to$ méthode numérique $\to$ outil du cours\} qui servira de \og carte routière\fg{} référencée dans l'introduction de chaque module suivant.

\noindent\textbf{Livrable prévu~:} tableau de correspondance (1 page) formulations/méthodes/outils, réutilisé en tête de chaque module 1 à 3.

\bigskip
\hrule
\bigskip

\chapter{Plateformes matériaux de la photonique intégrée~: SOI, SiN, InP}
\label{chap:m0-materiaux}
\etatunite{PLAN}

\noindent\textbf{Niveau~:} débutant (dans le contexte industriel) \hfill
\textbf{Position~:} Module 0, Chapitre 3 sur 4

\begin{objectifs}
\begin{itemize}[leftmargin=1.4em]
  \item \textbf{[Théorique]} Comparer les contrastes d'indice, fenêtres de transparence et non-linéarités des plateformes SOI (silicium sur isolant), SiN (nitrure de silicium) et InP (semi-conducteurs III-V actifs).
  \item \textbf{[Théorique]} Relier le choix de plateforme aux contraintes physiques déjà connues (dispersion, biréfringence, seuils de dommage/non-linéarité) et à des compromis applicatifs (pertes vs intégration active).
  \item \textbf{[Pratique]} Lire une fiche technique de plateforme (\emph{stack} de couches, épaisseurs, indices) telle que fournie typiquement par une fonderie, et en extraire les paramètres nécessaires à une première simulation.
  \item \textbf{[Pratique]} Construire en Python une table de dispersion matériau ($n(\lambda)$, modèle de Sellmeier) pour Si, SiN et SiO$_2$, réutilisable dans tous les modules suivants.
\end{itemize}
\end{objectifs}

\begin{prerequis}
\textbf{Théoriques~:} \cref{chap:m0-optique,chap:m0-em}; notions de dispersion chromatique et de non-linéarité optique — acquis de thèse pour la partie physique, à raccrocher au vocabulaire fonderie.\\
\textbf{Outils/logiciels~:} Python/NumPy/SciPy (modèles de Sellmeier).
\end{prerequis}

\noindent\textbf{Outils principaux~:} Python/SciPy (open source) pour les modèles de dispersion ; mention des bases de données d'indices usuelles (refractiveindex.info) comme ressource externe.

\noindent\textbf{Rappel de mise à niveau prévu~:} oui, court rappel sur les modèles de Sellmeier/Cauchy si non pratiqués depuis la thèse, formulé comme extension directe des notions de dispersion de matériaux laser déjà connues.

\noindent\textbf{Aperçu du contenu~:} présentation comparative structurée (tableau) SOI/SiN/InP~: contraste d'indice, pertes typiques, budget thermique, disponibilité fonderie (MPW), compatibilité CMOS — première introduction explicite du vocabulaire fonderie qui sera approfondi au Module~6.

\noindent\textbf{Livrable prévu~:} module Python \texttt{materiaux.py} (fonctions de dispersion Si/SiN/SiO$_2$) importé par les scripts de simulation des Modules 1 à 4.

\bigskip
\hrule
\bigskip

\chapter{Panorama des méthodes numériques et de l'écosystème logiciel}
\label{chap:m0-panorama}
\etatunite{PLAN}

\noindent\textbf{Niveau~:} débutant (dans le contexte industriel) \hfill
\textbf{Position~:} Module 0, Chapitre 4 sur 4

\begin{objectifs}
\begin{itemize}[leftmargin=1.4em]
  \item \textbf{[Théorique]} Positionner FDTD, EME, FEM et BPM (\emph{Beam Propagation Method}) les unes par rapport aux autres selon le type de problème résolu (propagation large bande vs modale, géométries invariantes vs quelconques).
  \item \textbf{[Pratique]} Cartographier l'écosystème logiciel du cours~: outils open source (Meep, scikit-fem/FEniCS, gdsfactory, KLayout) et outils commerciaux (Lumerical FDTD/MODE/INTERCONNECT/DEVICE, IPKISS, RSoft), avec statut de licence.
  \item \textbf{[Pratique]} Installer et valider un environnement Python de base commun à tout le cours (gestion d'environnement virtuel, dépendances scientifiques).
  \item \textbf{[Pratique]} Anticiper les besoins de licence académique/essai pour les outils commerciaux mobilisés dès le Module~4.
\end{itemize}
\end{objectifs}

\begin{prerequis}
\textbf{Théoriques~:} \cref{chap:m0-em} (tableau de correspondance formulations/méthodes).\\
\textbf{Outils/logiciels~:} Python $\geq$ 3.10, gestionnaire d'environnement (conda ou venv) ; à ce stade, aucune licence commerciale requise.
\end{prerequis}

\noindent\textbf{Outils principaux~:} vue d'ensemble uniquement — Meep (open source, MIT), scikit-fem/FEniCS (open source), gdsfactory/KLayout (open source), Lumerical/IPKISS/RSoft (commerciaux, licence académique ou essai à demander en amont du Module~4).

\noindent\textbf{Rappel de mise à niveau prévu~:} non.

\noindent\textbf{Aperçu du contenu~:} tableau de décision \og quelle méthode pour quel problème\fg{} et checklist d'installation d'environnement, conçus pour être référencés (pas redupliqués) en tête de chaque module ultérieur.

\noindent\textbf{Livrable prévu~:} environnement Python reproductible (\texttt{environment.yml} ou \texttt{requirements.txt}) validé par un script de test d'import de toutes les bibliothèques open source du cours.


% ============================================================
% MODULE 1 — FDTD open source
% ============================================================
\part{FDTD open source~: principes, Meep, validation vs théorie}
\label{part:module1}
% ============================================================
% Module 1 — FDTD open source (Meep)
% ÉTAT : PLAN.
% ============================================================

\chapter{Principes de la méthode FDTD~: schéma de Yee, stabilité, PML}
\label{chap:m1-principes}
\etatunite{PLAN}

\noindent\textbf{Niveau~:} débutant \hfill \textbf{Position~:} Module 1, Chapitre 1 sur 5

\begin{objectifs}
\begin{itemize}[leftmargin=1.4em]
  \item \textbf{[Théorique]} Dériver le schéma de discrétisation de Yee (grille décalée E/H) à partir des équations de Maxwell temporelles présentées en \cref{chap:m0-em}.
  \item \textbf{[Théorique]} Établir la condition de stabilité de Courant-Friedrichs-Lewy (CFL) et en déduire l'impact du pas spatial sur le pas temporel.
  \item \textbf{[Théorique]} Justifier physiquement le rôle des couches PML (\emph{Perfectly Matched Layer}) en s'appuyant sur le rappel d'apodisation introduit en \cref{chap:m0-em}.
  \item \textbf{[Pratique]} Traduire un choix de résolution spatiale (mailles/longueur d'onde) en coût mémoire et temps de calcul estimés, à la main, avant toute exécution.
\end{itemize}
\end{objectifs}

\begin{prerequis}
\textbf{Théoriques~:} \cref{chap:m0-em} (formulation temporelle de Maxwell, PML); \cref{chap:m0-optique} (modes guidés).\\
\textbf{Outils/logiciels~:} aucun encore ; chapitre principalement analytique/numérique \og à la main\fg{}.
\end{prerequis}

\noindent\textbf{Outils principaux~:} aucun solveur ; calculs de dimensionnement en Python/NumPy.

\noindent\textbf{Rappel de mise à niveau prévu~:} non (le rappel PML a déjà été placé en \cref{chap:m0-em}).

\noindent\textbf{Aperçu du contenu~:} construction pas à pas du schéma de Yee 1D puis 2D, analyse de dispersion numérique (différence entre vitesse de phase discrète et continue) — point de vigilance repris dans l'exercice diagnostic du \cref{chap:m1-guide}.

\noindent\textbf{Livrable prévu~:} note de dimensionnement (résolution, pas de temps, épaisseur PML) pour le cas guide droit qui sera simulé au \cref{chap:m1-guide}.

\bigskip\hrule\bigskip

\chapter{Prise en main de Meep~: installation, scripts Python, unités normalisées}
\label{chap:m1-meep}
\etatunite{PLAN}

\noindent\textbf{Niveau~:} débutant \hfill \textbf{Position~:} Module 1, Chapitre 2 sur 5

\begin{objectifs}
\begin{itemize}[leftmargin=1.4em]
  \item \textbf{[Pratique]} Installer Meep (via conda-forge) et valider l'installation par un script minimal.
  \item \textbf{[Théorique]} Convertir entre unités physiques (SI) et unités normalisées de Meep ($a$, $c=1$), source classique d'erreur pour un nouvel utilisateur.
  \item \textbf{[Pratique]} Construire les objets de base d'une simulation Meep~: \texttt{Simulation}, \texttt{Medium}, \texttt{Source}, \texttt{Monitor/Volume}.
  \item \textbf{[Pratique]} Exécuter une première simulation \og jouet\fg{} (résonateur 2D) et visualiser les champs avec Matplotlib.
\end{itemize}
\end{objectifs}

\begin{prerequis}
\textbf{Théoriques~:} \cref{chap:m1-principes}.\\
\textbf{Outils/logiciels~:} Meep $\geq$ 1.27 (conda-forge), h5py, Matplotlib.
\end{prerequis}

\noindent\textbf{Outils principaux~:} Meep (open source, GPL).

\noindent\textbf{Rappel de mise à niveau prévu~:} oui, court encart sur le système d'unités normalisées (analogie avec les unités \og naturelles\fg{} parfois utilisées en physique laser/optique non linéaire).

\noindent\textbf{Aperçu du contenu~:} tutoriel guidé d'installation et de première simulation, calé sur l'exemple canonique Meep (résonateur) mais reformulé pour l'objectif final du module (guide droit).

\noindent\textbf{Livrable prévu~:} environnement Meep fonctionnel + script \texttt{hello\_meep.py} exécuté avec succès (figure de champ en sortie).

\bigskip\hrule\bigskip

\chapter{Simulation d'un guide droit~: validation de $n_{\text{eff}}$ et des pertes}
\label{chap:m1-guide}
\etatunite{PLAN}

\noindent\textbf{Niveau~:} intermédiaire \hfill \textbf{Position~:} Module 1, Chapitre 3 sur 5

\begin{objectifs}
\begin{itemize}[leftmargin=1.4em]
  \item \textbf{[Pratique]} Construire en Meep un guide droit rectangulaire (plateforme SOI, cf. \cref{chap:m0-materiaux}) avec source modale et moniteurs de flux.
  \item \textbf{[Théorique]} Extraire $n_{\text{eff}}$ à partir d'une simulation FDTD (analyse de phase ou comparaison à la solution analytique du \cref{chap:m0-optique}) et discuter les sources d'écart.
  \item \textbf{[Pratique]} Quantifier les pertes de propagation numériques et les distinguer des pertes physiques (rugosité non modélisée à ce stade).
  \item \textbf{[Théorique]} Diagnostiquer un écart de résultat lié à un sous-échantillonnage spatial (lien direct avec la dispersion numérique du \cref{chap:m1-principes}).
\end{itemize}
\end{objectifs}

\begin{prerequis}
\textbf{Théoriques~:} \cref{chap:m1-principes,chap:m0-optique,chap:m0-materiaux}.\\
\textbf{Outils/logiciels~:} Meep, script \texttt{materiaux.py} du \cref{chap:m0-materiaux}.
\end{prerequis}

\noindent\textbf{Outils principaux~:} Meep.

\noindent\textbf{Rappel de mise à niveau prévu~:} non.

\noindent\textbf{Aperçu du contenu~:} ce chapitre fournit le cas de référence \og guide droit\fg{} qui sera repris à l'identique en Lumerical FDTD au \cref{chap:m4-fdtd-lumerical}, pour une comparaison directe outil open source / outil commercial.

\noindent\textbf{Livrable prévu~:} script \texttt{guide\_droit\_meep.py} + figure $n_{\text{eff}}$ simulé vs analytique + tableau d'écart relatif, réutilisés comme référence de validation au Module~4.

\bigskip\hrule\bigskip

\chapter{Coupleurs directionnels~: simulation et étude paramétrique}
\label{chap:m1-coupleur}
\etatunite{PLAN}

\noindent\textbf{Niveau~:} intermédiaire \hfill \textbf{Position~:} Module 1, Chapitre 4 sur 5

\begin{objectifs}
\begin{itemize}[leftmargin=1.4em]
  \item \textbf{[Théorique]} Relier le couplage entre guides voisins à la théorie des modes couplés, en s'appuyant sur l'analogie avec le couplage entre cavités/modes de fibres jumelles.
  \item \textbf{[Pratique]} Construire un coupleur directionnel en Meep et balayer la longueur de couplage pour retrouver la longueur de battement $L_\pi$.
  \item \textbf{[Pratique]} Automatiser un balayage paramétrique (gap, longueur) en Python et tracer le taux de couplage résultant.
  \item \textbf{[Théorique]} Relier la sensibilité du couplage au gap à une première contrainte de fabrication (tolérance de lithographie), annonçant le \cref{chap:m1-postraitement}.
\end{itemize}
\end{objectifs}

\begin{prerequis}
\textbf{Théoriques~:} \cref{chap:m1-guide}; théorie des modes couplés — notion à rappeler brièvement si non pratiquée depuis la thèse.\\
\textbf{Outils/logiciels~:} Meep, NumPy pour le post-traitement du balayage.
\end{prerequis}

\noindent\textbf{Outils principaux~:} Meep.

\noindent\textbf{Rappel de mise à niveau prévu~:} oui, court encart sur la théorie des modes couplés (CMT), formulé comme extension de la théorie des cavités couplées déjà connue.

\noindent\textbf{Aperçu du contenu~:} étude paramétrique automatisée servant de gabarit méthodologique repris tel quel avec l'API \texttt{lumapi} au \cref{chap:m4-api}.

\noindent\textbf{Livrable prévu~:} courbe taux de couplage vs longueur/gap + valeur de $L_\pi$ extraite, avec code de balayage réutilisable.

\bigskip\hrule\bigskip

\chapter{Sources, moniteurs et post-traitement~: S-paramètres et spectres}
\label{chap:m1-postraitement}
\etatunite{PLAN}

\noindent\textbf{Niveau~:} intermédiaire \hfill \textbf{Position~:} Module 1, Chapitre 5 sur 5

\begin{objectifs}
\begin{itemize}[leftmargin=1.4em]
  \item \textbf{[Théorique]} Définir les S-paramètres d'un composant photonique par analogie avec les paramètres de transmission/réflexion déjà connus en optique (matrices ABCD, coefficients de Fresnel).
  \item \textbf{[Pratique]} Extraire un spectre large bande en une seule simulation FDTD (source impulsionnelle + transformée de Fourier des moniteurs de flux).
  \item \textbf{[Pratique]} Normaliser correctement un spectre (simulation de référence sans dispositif) — erreur fréquente de débutant.
  \item \textbf{[Pratique]} Produire une figure de spectre publiable (échelle dB, légende, bande utile) qui servira de format standard pour tout le reste du cours.
\end{itemize}
\end{objectifs}

\begin{prerequis}
\textbf{Théoriques~:} \cref{chap:m1-coupleur}; notion de coefficients de transmission/réflexion — acquis de thèse.\\
\textbf{Outils/logiciels~:} Meep, h5py, Matplotlib.
\end{prerequis}

\noindent\textbf{Outils principaux~:} Meep.

\noindent\textbf{Rappel de mise à niveau prévu~:} non.

\noindent\textbf{Aperçu du contenu~:} ce chapitre fixe le format standard de figure spectrale (dB, bande utile, normalisation) réutilisé sans redéfinition dans les Modules 2 à 4 ; clôt le Module~1 sur le cas coupleur directionnel entièrement caractérisé.

\noindent\textbf{Livrable prévu~:} fonction Python \texttt{plot\_spectrum()} réutilisable + spectre de transmission normalisé du coupleur directionnel du \cref{chap:m1-coupleur}, clôturant le portfolio du Module~1.


% ============================================================
% MODULE 2 — EME
% ============================================================
\part{EME (Eigenmode Expansion)~: principes et cas d'usage vs FDTD}
\label{part:module2}
% ============================================================
% Module 2 — EME (Eigenmode Expansion)
% ÉTAT : PLAN.
% ============================================================

\chapter{Principes de l'Eigenmode Expansion~: décomposition modale et matrices S}
\label{chap:m2-principes}
\etatunite{PLAN}

\noindent\textbf{Niveau~:} intermédiaire \hfill \textbf{Position~:} Module 2, Chapitre 1 sur 3

\begin{objectifs}
\begin{itemize}[leftmargin=1.4em]
  \item \textbf{[Théorique]} Formuler la propagation d'un champ comme une somme pondérée de modes propres locaux, en s'appuyant explicitement sur l'analogie avec la décomposition en modes de cavité laser déjà maîtrisée.
  \item \textbf{[Théorique]} Construire la matrice de raccordement (\emph{mode matching}) entre deux sections de guide d'indices effectifs différents et en déduire la matrice S d'une structure segmentée.
  \item \textbf{[Théorique]} Identifier les hypothèses de validité de l'EME (invariance longitudinale par section, base modale complète) et leurs limites.
  \item \textbf{[Pratique]} Calculer \og à la main\fg{}, sur un cas à deux sections, le coefficient de transmission par la méthode des matrices S, avant toute utilisation d'un solveur.
\end{itemize}
\end{objectifs}

\begin{prerequis}
\textbf{Théoriques~:} \cref{chap:m0-optique} (modes guidés); \cref{chap:m0-em} (tableau de correspondance formulations/méthodes); matrices ABCD de l'optique gaussienne — acquis de thèse, à transposer.\\
\textbf{Outils/logiciels~:} Python/NumPy (algèbre matricielle) uniquement à ce stade.
\end{prerequis}

\noindent\textbf{Outils principaux~:} aucun solveur dédié — construction manuelle des matrices en Python pour ancrer la théorie avant l'outillage du \cref{chap:m2-vs-fdtd}.

\noindent\textbf{Rappel de mise à niveau prévu~:} oui, encart reliant explicitement les matrices S de l'EME aux matrices ABCD de l'optique gaussienne déjà connues, comme changement de base conceptuel plutôt que nouveauté.

\noindent\textbf{Aperçu du contenu~:} construction progressive~: un mode $\to$ deux modes $\to$ discontinuité $\to$ cascade de sections, avant généralisation aux solveurs EME du \cref{chap:m2-application}.

\noindent\textbf{Livrable prévu~:} script Python de calcul de matrice S par recouvrement modal sur un cas à 2 sections (discontinuité de largeur de guide), avec comparaison à un cas limite analytique connu.

\bigskip\hrule\bigskip

\chapter{EME vs FDTD~: critères de choix et limites}
\label{chap:m2-vs-fdtd}
\etatunite{PLAN}

\noindent\textbf{Niveau~:} intermédiaire \hfill \textbf{Position~:} Module 2, Chapitre 2 sur 3

\begin{objectifs}
\begin{itemize}[leftmargin=1.4em]
  \item \textbf{[Théorique]} Établir une grille de décision EME vs FDTD selon le type de structure (invariante par section vs géométrie quelconque, longueur du dispositif, besoin de réflexions multiples).
  \item \textbf{[Pratique]} Comparer, sur le \emph{même} guide droit déjà validé au \cref{chap:m1-guide}, le temps de calcul et la précision obtenus en EME vs FDTD.
  \item \textbf{[Théorique]} Expliquer pourquoi l'EME est particulièrement adaptée aux structures longues et progressives (tapers) là où FDTD devient prohibitif en temps de calcul.
  \item \textbf{[Pratique]} Diagnostiquer un cas où l'EME donne un résultat incorrect faute de nombre de modes suffisant dans la base (troncature modale).
\end{itemize}
\end{objectifs}

\begin{prerequis}
\textbf{Théoriques~:} \cref{chap:m2-principes,chap:m1-guide}.\\
\textbf{Outils/logiciels~:} solveur EME open source retenu pour le cours (à confirmer~: EMEpy ou combinaison solveur modal scikit-fem + propagation matricielle maison).
\end{prerequis}

\noindent\textbf{Outils principaux~:} solveur EME open source (Python) — choix d'outil à trancher en tête de chapitre selon maturité/maintenance au moment de la rédaction.

\noindent\textbf{Rappel de mise à niveau prévu~:} non.

\noindent\textbf{Aperçu du contenu~:} tableau de décision méthodologique + benchmark chiffré temps/précision sur le cas guide droit, complétant le tableau \og quelle méthode pour quel problème\fg{} du \cref{chap:m0-panorama}.

\noindent\textbf{Livrable prévu~:} tableau comparatif EME/FDTD (temps de calcul, précision, cas d'usage recommandé) sur le cas guide droit de référence.

\bigskip\hrule\bigskip

\chapter{Application EME~: tapers et transitions adiabatiques}
\label{chap:m2-application}
\etatunite{PLAN}

\noindent\textbf{Niveau~:} intermédiaire \hfill \textbf{Position~:} Module 2, Chapitre 3 sur 3

\begin{objectifs}
\begin{itemize}[leftmargin=1.4em]
  \item \textbf{[Théorique]} Définir le critère d'adiabaticité d'une transition de guide (taper) et le relier à la notion de croisement de modes évité.
  \item \textbf{[Pratique]} Simuler en EME un taper linéaire puis un taper optimisé (parabolique) et comparer leurs pertes de conversion modale.
  \item \textbf{[Pratique]} Réaliser un balayage de longueur de taper pour identifier la longueur minimale garantissant une conversion adiabatique sous un seuil de perte donné.
  \item \textbf{[Théorique]} Relier la longueur de taper obtenue à une contrainte fonderie réaliste (encombrement de puce, budget de longueur imposé par un PDK).
\end{itemize}
\end{objectifs}

\begin{prerequis}
\textbf{Théoriques~:} \cref{chap:m2-vs-fdtd}.\\
\textbf{Outils/logiciels~:} solveur EME retenu (cf. \cref{chap:m2-vs-fdtd}).
\end{prerequis}

\noindent\textbf{Outils principaux~:} solveur EME open source.

\noindent\textbf{Rappel de mise à niveau prévu~:} non.

\noindent\textbf{Aperçu du contenu~:} premier chapitre du cours à faire explicitement le lien avec une contrainte de PDK (\S \emph{Lien avec le flot fonderie}), en préparation du Module~5.

\begin{flotfonderie}
Aperçu (à détailler en rédaction complète)~: budget de longueur de taper imposé par un PDK type (ex. $\leq$~\SI{20}{\micro\metre} pour une transition ridge-to-strip) et son impact direct sur le choix de profil (linéaire vs parabolique) simulé dans ce chapitre.
\end{flotfonderie}

\noindent\textbf{Livrable prévu~:} courbe perte de conversion vs longueur de taper (linéaire et parabolique) + recommandation de longueur minimale sous contrainte fonderie donnée.


% ============================================================
% MODULE 3 — FEM
% ============================================================
\part{FEM~: principes, cas d'usage et outils associés}
\label{part:module3}
% ============================================================
% Module 3 — FEM (Finite Element Method)
% ÉTAT : PLAN.
% ============================================================

\chapter{Principes de la méthode des éléments finis pour les modes de guide}
\label{chap:m3-principes}
\etatunite{PLAN}

\noindent\textbf{Niveau~:} intermédiaire \hfill \textbf{Position~:} Module 3, Chapitre 1 sur 4

\begin{objectifs}
\begin{itemize}[leftmargin=1.4em]
  \item \textbf{[Théorique]} Dériver la formulation faible (variationnelle) de l'équation d'onde vectorielle pour le calcul de modes, à partir de la formulation forte déjà rencontrée en \cref{chap:m0-em}.
  \item \textbf{[Théorique]} Expliquer le principe de discrétisation par éléments (fonctions de base locales) et le comparer à la discrétisation par grille régulière de la FDTD (\cref{chap:m1-principes}).
  \item \textbf{[Théorique]} Justifier pourquoi la FEM est la méthode de référence pour des géométries non rectangulaires (guides gravés en biseau, coins arrondis) là où FDTD/EME sont moins naturelles.
  \item \textbf{[Pratique]} Résoudre \og à la main\fg{} un problème aux valeurs propres 1D simplifié (discrétisation à 3-5 éléments) pour retrouver le principe avant le passage à un outil dédié.
\end{itemize}
\end{objectifs}

\begin{prerequis}
\textbf{Théoriques~:} \cref{chap:m0-em,chap:m1-principes}; algèbre linéaire (problèmes aux valeurs propres généralisés) — acquis de thèse.\\
\textbf{Outils/logiciels~:} Python/NumPy/SciPy (\texttt{scipy.sparse.linalg.eigs}) pour l'illustration \og à la main\fg{}.
\end{prerequis}

\noindent\textbf{Outils principaux~:} aucun solveur FEM dédié à ce stade — construction pédagogique minimale en Python/SciPy.

\noindent\textbf{Rappel de mise à niveau prévu~:} oui, court rappel sur les problèmes aux valeurs propres généralisés $A x = \lambda B x$, si rouillé, avant leur réapparition sous forme de calcul modal FEM.

\noindent\textbf{Aperçu du contenu~:} parallèle explicite formulation forte (FDTD) / formulation faible (FEM), consolidant le tableau de correspondance du \cref{chap:m0-panorama}.

\noindent\textbf{Livrable prévu~:} notebook Python de résolution d'un problème aux valeurs propres 1D minimal, servant de \og pont\fg{} conceptuel avant l'utilisation d'un solveur FEM complet au \cref{chap:m3-outils}.

\bigskip\hrule\bigskip

\chapter{Maillage et convergence~: stratégies pour les guides photoniques}
\label{chap:m3-maillage}
\etatunite{PLAN}

\noindent\textbf{Niveau~:} intermédiaire \hfill \textbf{Position~:} Module 3, Chapitre 2 sur 4

\begin{objectifs}
\begin{itemize}[leftmargin=1.4em]
  \item \textbf{[Théorique]} Définir les critères de qualité d'un maillage (raffinement aux interfaces d'indice, ratio d'aspect des éléments) et leur impact sur la précision du $n_{\text{eff}}$ calculé.
  \item \textbf{[Pratique]} Réaliser une étude de convergence en maillage (raffinement progressif) et déterminer un critère d'arrêt quantitatif.
  \item \textbf{[Théorique]} Relier une non-convergence à une cause physique/numérique précise (maillage insuffisant à l'interface cœur/gaine, coin non résolu) — première préparation à l'exercice diagnostic du module.
\end{itemize}
\end{objectifs}

\begin{prerequis}
\textbf{Théoriques~:} \cref{chap:m3-principes}.\\
\textbf{Outils/logiciels~:} outil de maillage associé au solveur FEM retenu (cf. \cref{chap:m3-outils}) — Gmsh pressenti comme générateur de maillage open source standard.
\end{prerequis}

\noindent\textbf{Outils principaux~:} Gmsh (maillage, open source) en préparation de l'utilisation effective au \cref{chap:m3-outils}.

\noindent\textbf{Rappel de mise à niveau prévu~:} non.

\noindent\textbf{Aperçu du contenu~:} méthodologie générique de convergence (transférable à tout solveur FEM commercial rencontré plus tard, y compris hors du cours), illustrée sur le guide droit de référence du \cref{chap:m1-guide}.

\noindent\textbf{Livrable prévu~:} courbe de convergence $n_{\text{eff}}$ vs finesse de maillage + critère de maillage retenu, documenté pour réutilisation dans tous les calculs FEM ultérieurs du cours.

\bigskip\hrule\bigskip

\chapter{Outils FEM open source pour le calcul modal}
\label{chap:m3-outils}
\etatunite{PLAN}

\noindent\textbf{Niveau~:} intermédiaire \hfill \textbf{Position~:} Module 3, Chapitre 3 sur 4

\begin{objectifs}
\begin{itemize}[leftmargin=1.4em]
  \item \textbf{[Pratique]} Installer et prendre en main un solveur FEM open source pour le calcul modal (scikit-fem ou FEniCSx~: choix à confirmer en rédaction selon maturité de l'exemple guide d'onde au moment de la rédaction).
  \item \textbf{[Pratique]} Reconstruire en FEM le guide droit de référence (\cref{chap:m1-guide}) et comparer $n_{\text{eff}}$ obtenu à FDTD (Meep) et à la solution analytique (\cref{chap:m0-optique}).
  \item \textbf{[Pratique]} Calculer un cas que FDTD/EME traitent mal nativement~: un guide à coin arrondi ou à flanc gravé en biseau, réaliste d'un point de vue fonderie.
  \item \textbf{[Théorique]} Interpréter un résultat de convergence à la lumière des critères établis au \cref{chap:m3-maillage}.
\end{itemize}
\end{objectifs}

\begin{prerequis}
\textbf{Théoriques~:} \cref{chap:m3-maillage,chap:m1-guide}.\\
\textbf{Outils/logiciels~:} solveur FEM open source retenu + Gmsh.
\end{prerequis}

\noindent\textbf{Outils principaux~:} solveur FEM open source (scikit-fem ou FEniCSx) + Gmsh, open source.

\noindent\textbf{Rappel de mise à niveau prévu~:} non.

\noindent\textbf{Aperçu du contenu~:} triangulation croisée des trois méthodes du cours (FDTD/EME/FEM) sur un même cas de référence, consolidant le benchmark méthodologique amorcé au \cref{chap:m2-vs-fdtd}.

\noindent\textbf{Livrable prévu~:} tableau triple comparatif $n_{\text{eff}}$ (analytique / FDTD / FEM) sur le guide droit + calcul FEM d'un guide à géométrie non idéale (coin arrondi).

\bigskip\hrule\bigskip

\chapter{Applications électro-optiques~: introduction aux couplages FEM multiphysiques}
\label{chap:m3-multiphysique}
\etatunite{PLAN}

\noindent\textbf{Niveau~:} avancé \hfill \textbf{Position~:} Module 3, Chapitre 4 sur 4

\begin{objectifs}
\begin{itemize}[leftmargin=1.4em]
  \item \textbf{[Théorique]} Présenter le principe d'un couplage multiphysique FEM (par exemple thermo-optique~: champ de température $\to$ variation d'indice $\to$ recalcul modal), sans viser l'exhaustivité.
  \item \textbf{[Théorique]} Situer ce type de couplage par rapport à un cas déjà connu du physicien~: effet thermo-optique dans un cristal laser, effet électro-optique dans un modulateur Pockels.
  \item \textbf{[Pratique]} Réaliser un exemple minimal de couplage thermo-optique (profil de température imposé $\to$ décalage de $n_{\text{eff}}$) sur le guide de référence.
  \item \textbf{[Théorique]} Anticiper le lien avec les composants actifs qui seront réutilisés en modèle compact à l'INTERCONNECT (\cref{chap:m4-interconnect}).
\end{itemize}
\end{objectifs}

\begin{prerequis}
\textbf{Théoriques~:} \cref{chap:m3-outils}; effet thermo-optique/électro-optique — acquis de thèse pour la physique, nouveauté pour la mise en œuvre numérique couplée.\\
\textbf{Outils/logiciels~:} solveur FEM retenu, éventuellement extension multiphysique si disponible dans l'outil open source choisi (sinon couplage manuel via script Python).
\end{prerequis}

\noindent\textbf{Outils principaux~:} solveur FEM open source (couplage manuel thermique/optique via script Python si l'outil ne propose pas de multiphysique intégrée).

\noindent\textbf{Rappel de mise à niveau prévu~:} oui, court rappel sur le coefficient thermo-optique $dn/dT$ et son ordre de grandeur pour Si/SiN, par analogie avec les effets thermiques déjà rencontrés en physique laser.

\noindent\textbf{Aperçu du contenu~:} chapitre volontairement introductif (non exhaustif), positionné comme porte d'entrée vers des couplages multiphysiques plus poussés hors du périmètre strict de ce cours.

\noindent\textbf{Livrable prévu~:} courbe de décalage de $n_{\text{eff}}$ en fonction d'un $\Delta T$ imposé, clôturant le Module~3 sur un cas connectable à un composant actif réel (modulateur thermo-optique).


% ============================================================
% MODULE 4 — Ansys Lumerical pour le PIC design
% ============================================================
\part{Ansys Lumerical pour le PIC design}
\label{part:module4}
% ============================================================
% Module 4 — Ansys Lumerical pour le PIC design
% ÉTAT : PLAN.
% Module autonome dédié à la suite Lumerical (outil quasi-standard
% de l'industrie), avec sa propre progression débutant -> avancé
% au sein de l'outil, indépendamment du niveau théorique déjà
% couvert par les Modules 1-3.
% ============================================================

\chapter{Prise en main de l'interface Lumerical}
\label{chap:m4-interface}
\etatunite{PLAN}

\noindent\textbf{Niveau~:} débutant (dans l'outil) \hfill \textbf{Position~:} Module 4, Chapitre 1 sur 8

\begin{objectifs}
\begin{itemize}[leftmargin=1.4em]
  \item \textbf{[Pratique]} Naviguer dans l'interface CAD Lumerical~: hiérarchie de simulation, objets de structure, objets de simulation, objets d'analyse.
  \item \textbf{[Pratique]} Distinguer les produits de la suite (FDTD Solutions, MODE Solutions, INTERCONNECT, DEVICE) et leur rôle respectif dans un flot de conception PIC.
  \item \textbf{[Pratique]} Identifier le modèle de licence (FlexLM, licences flottantes/nommées) et anticiper une demande de licence académique ou d'essai.
  \item \textbf{[Théorique]} Relier chaque produit Lumerical à la méthode numérique sous-jacente déjà vue (FDTD Solutions $\leftrightarrow$ Module 1, MODE/EME $\leftrightarrow$ Module 2, DEVICE $\leftrightarrow$ FEM du Module 3).
\end{itemize}
\end{objectifs}

\begin{prerequis}
\textbf{Théoriques~:} \cref{chap:m0-panorama} (tableau outils/méthodes).\\
\textbf{Outils/logiciels~:} Ansys Lumerical (licence académique ou essai à obtenir avant ce chapitre — délai à anticiper, cf. remarque de planification en tête de module).
\end{prerequis}

\noindent\textbf{Outils principaux~:} Ansys Lumerical (commercial ; licence académique/essai nécessaire).

\noindent\textbf{Rappel de mise à niveau prévu~:} non.

\noindent\textbf{Aperçu du contenu~:} visite guidée de l'interface avec vocabulaire anglais standard (\emph{simulation region}, \emph{mesh override}, \emph{monitor}), posant le socle commun aux \cref{chap:m4-fdtd-lumerical,chap:m4-mode}.

\noindent\textbf{Livrable prévu~:} note de configuration de l'environnement Lumerical (licence, version, chemin d'installation) + capture d'écran annotée de la hiérarchie de simulation.

\bigskip\hrule\bigskip

\chapter{FDTD Solutions~: reconstruire le guide droit et le coupleur directionnel}
\label{chap:m4-fdtd-lumerical}
\etatunite{PLAN}

\noindent\textbf{Niveau~:} débutant $\to$ intermédiaire \hfill \textbf{Position~:} Module 4, Chapitre 2 sur 8

\begin{objectifs}
\begin{itemize}[leftmargin=1.4em]
  \item \textbf{[Pratique]} Reconstruire dans FDTD Solutions le guide droit du \cref{chap:m1-guide} et le coupleur directionnel du \cref{chap:m1-coupleur}.
  \item \textbf{[Pratique]} Comparer directement $n_{\text{eff}}$, pertes et taux de couplage obtenus en Lumerical FDTD à ceux obtenus en Meep, en quantifiant les écarts.
  \item \textbf{[Théorique]} Expliquer les différences de paramétrage entre Meep et FDTD Solutions (maillage automatique vs manuel, \emph{mesh override regions}, PML avancées) pour un même schéma numérique sous-jacent.
  \item \textbf{[Pratique]} Identifier, sur un résultat divergent, si l'écart provient du maillage, de la définition matériau, ou d'une différence de convention (première mise en pratique du réflexe diagnostic).
\end{itemize}
\end{objectifs}

\begin{prerequis}
\textbf{Théoriques~:} \cref{chap:m1-guide,chap:m1-coupleur}; \cref{chap:m4-interface}.\\
\textbf{Outils/logiciels~:} Lumerical FDTD Solutions.
\end{prerequis}

\noindent\textbf{Outils principaux~:} Lumerical FDTD Solutions (commercial) ; équivalent open source déjà vu~: Meep (\cref{chap:m1-meep}).

\noindent\textbf{Rappel de mise à niveau prévu~:} non.

\noindent\textbf{Aperçu du contenu~:} chapitre de \og transfert de compétence\fg{} explicite Meep $\to$ Lumerical, sur des cas déjà maîtrisés, pour isoler l'apprentissage de l'outil de l'apprentissage de la physique.

\noindent\textbf{Livrable prévu~:} tableau comparatif Meep vs Lumerical FDTD ($n_{\text{eff}}$, pertes, $L_\pi$) sur les deux cas de référence du Module~1.

\bigskip\hrule\bigskip

\chapter{MODE Solutions / solveur EME~: analyse modale, dispersion, tapers}
\label{chap:m4-mode}
\etatunite{PLAN}

\noindent\textbf{Niveau~:} intermédiaire \hfill \textbf{Position~:} Module 4, Chapitre 3 sur 8

\begin{objectifs}
\begin{itemize}[leftmargin=1.4em]
  \item \textbf{[Pratique]} Utiliser le solveur modal FDE (\emph{Finite Difference Eigenmode}) de MODE Solutions et le solveur EME associé.
  \item \textbf{[Pratique]} Calculer la courbe de dispersion $n_{\text{eff}}(\lambda)$ d'un guide et la comparer à celle obtenue analytiquement (\cref{chap:m0-optique}) et en FEM (\cref{chap:m3-outils}).
  \item \textbf{[Pratique]} Reconstruire le taper adiabatique du \cref{chap:m2-application} dans le solveur EME de MODE Solutions et comparer les pertes de conversion.
  \item \textbf{[Théorique]} Expliquer pourquoi MODE Solutions utilise un solveur FDE (et non FEM) pour l'analyse modale, en revenant sur le tableau de correspondance du \cref{chap:m0-panorama}.
\end{itemize}
\end{objectifs}

\begin{prerequis}
\textbf{Théoriques~:} \cref{chap:m2-application,chap:m3-outils}; \cref{chap:m4-interface}.\\
\textbf{Outils/logiciels~:} Lumerical MODE Solutions.
\end{prerequis}

\noindent\textbf{Outils principaux~:} Lumerical MODE Solutions (commercial) ; équivalent open source déjà vu~: solveur EME du Module~2 et solveur FEM du Module~3.

\noindent\textbf{Rappel de mise à niveau prévu~:} non.

\noindent\textbf{Aperçu du contenu~:} triangulation à trois termes (EME open source / FEM open source / MODE Solutions) sur le cas taper, consolidant la confiance de l'étudiant dans les résultats commerciaux via des résultats déjà validés en open source.

\noindent\textbf{Livrable prévu~:} courbe de dispersion + tableau comparatif pertes de taper (EME open source vs MODE Solutions).

\bigskip\hrule\bigskip

\chapter{Automatisation avec l'API Python Lumerical (\texttt{lumapi})}
\label{chap:m4-api}
\etatunite{PLAN}

\noindent\textbf{Niveau~:} intermédiaire \hfill \textbf{Position~:} Module 4, Chapitre 4 sur 8

\begin{objectifs}
\begin{itemize}[leftmargin=1.4em]
  \item \textbf{[Pratique]} Piloter Lumerical FDTD/MODE depuis Python via \texttt{lumapi} au lieu de l'interface graphique.
  \item \textbf{[Pratique]} Reproduire par script le balayage paramétrique gap/longueur du coupleur directionnel (\cref{chap:m1-coupleur}), cette fois en Lumerical, en réutilisant le même gabarit méthodologique.
  \item \textbf{[Pratique]} Extraire automatiquement des résultats (S-paramètres, champs) vers des structures NumPy pour post-traitement unifié avec le reste du cours.
  \item \textbf{[Théorique]} Discuter l'intérêt de l'automatisation pour l'optimisation et l'analyse de tolérance, en préparation du \cref{chap:m4-tolerance}.
\end{itemize}
\end{objectifs}

\begin{prerequis}
\textbf{Théoriques~:} \cref{chap:m1-coupleur,chap:m4-fdtd-lumerical}.\\
\textbf{Outils/logiciels~:} Lumerical + module Python \texttt{lumapi} (fourni avec l'installation Lumerical), Python $\geq$ 3.10.
\end{prerequis}

\noindent\textbf{Outils principaux~:} Lumerical \texttt{lumapi} (commercial, API Python).

\noindent\textbf{Rappel de mise à niveau prévu~:} non.

\noindent\textbf{Aperçu du contenu~:} bascule du \og tout interface graphique\fg{} vers un flot scriptable, condition nécessaire aux chapitres avancés du module (tolérance, PDK).

\noindent\textbf{Livrable prévu~:} script \texttt{lumapi} de balayage automatisé (gap/longueur) + jeu de données exporté, réutilisé au \cref{chap:m4-tolerance}.

\bigskip\hrule\bigskip

\chapter{INTERCONNECT~: simulation circuit-level et assemblage de composants}
\label{chap:m4-interconnect}
\etatunite{PLAN}

\noindent\textbf{Niveau~:} intermédiaire $\to$ avancé \hfill \textbf{Position~:} Module 4, Chapitre 5 sur 8

\begin{objectifs}
\begin{itemize}[leftmargin=1.4em]
  \item \textbf{[Théorique]} Distinguer simulation \emph{device-level} (champ électromagnétique complet, FDTD/MODE) et simulation \emph{circuit-level} (composants réduits à des modèles compacts/S-paramètres) — analogie avec le passage optique physique $\to$ schéma optique en physique laser.
  \item \textbf{[Pratique]} Assembler dans INTERCONNECT un circuit simple (interféromètre de Mach-Zehnder) à partir des S-paramètres extraits en FDTD/MODE dans les chapitres précédents.
  \item \textbf{[Pratique]} Comparer la réponse spectrale obtenue circuit-level à une simulation device-level complète (quand elle est encore tractable) pour valider la méthode.
  \item \textbf{[Théorique]} Justifier pourquoi le circuit-level devient indispensable dès que le nombre de composants augmente (limite pratique du device-level).
\end{itemize}
\end{objectifs}

\begin{prerequis}
\textbf{Théoriques~:} \cref{chap:m4-fdtd-lumerical,chap:m4-mode,chap:m1-postraitement} (format standard de spectre).\\
\textbf{Outils/logiciels~:} Lumerical INTERCONNECT.
\end{prerequis}

\noindent\textbf{Outils principaux~:} Lumerical INTERCONNECT (commercial).

\noindent\textbf{Rappel de mise à niveau prévu~:} non.

\noindent\textbf{Aperçu du contenu~:} premier exemple complet de \emph{co-design flow} device $\to$ circuit du cours, préfigurant le mini-projet du \cref{chap:m4-projet}.

\noindent\textbf{Livrable prévu~:} circuit INTERCONNECT d'un MZI simple + spectre de transmission comparé device-level vs circuit-level.

\bigskip\hrule\bigskip

\chapter{Import/export de PDK fonderie dans Lumerical}
\label{chap:m4-pdk}
\etatunite{PLAN}

\noindent\textbf{Niveau~:} avancé \hfill \textbf{Position~:} Module 4, Chapitre 6 sur 8

\begin{objectifs}
\begin{itemize}[leftmargin=1.4em]
  \item \textbf{[Pratique]} Importer un PDK photonique (open ou académique, ex. PDK d'une fonderie MPW type AIM Photonics/imec/CORNERSTONE selon disponibilité) dans l'environnement Lumerical.
  \item \textbf{[Théorique]} Distinguer composant \og fait maison\fg{} (simulé de A à Z, Modules 1-3) et composant \og PDK\fg{} (pré-caractérisé, modèle compact fourni par la fonderie).
  \item \textbf{[Pratique]} Utiliser un composant pré-caractérisé du PDK dans un circuit INTERCONNECT (\cref{chap:m4-interconnect}) sans le resimuler intégralement.
  \item \textbf{[Théorique]} Identifier les contraintes de design imposées par le PDK (couches disponibles, largeurs/gaps minimaux) et leur impact sur les choix déjà faits dans le cours (ex. gap du coupleur du \cref{chap:m1-coupleur}).
\end{itemize}
\end{objectifs}

\begin{prerequis}
\textbf{Théoriques~:} \cref{chap:m4-interconnect}; \cref{chap:m0-materiaux} (vocabulaire fonderie).\\
\textbf{Outils/logiciels~:} Lumerical + kit PDK compatible (accès à confirmer selon disponibilité académique).
\end{prerequis}

\noindent\textbf{Outils principaux~:} Lumerical + PDK fonderie (accès nécessitant généralement un accord de confidentialité/licence académique, à anticiper).

\noindent\textbf{Rappel de mise à niveau prévu~:} oui, encart de vocabulaire fonderie (PDK, \emph{compact model}, \emph{process design kit}) formulé pour un public qui découvre ce jargon pour la première fois.

\noindent\textbf{Aperçu du contenu~:} bascule explicite du \og monde simulation académique\fg{} vers le \og monde fonderie industrielle\fg{}, articulée avec le Module~5 (layout/PDK) et le Module~6 (interaction fonderie).

\begin{flotfonderie}
Aperçu (à détailler en rédaction complète)~: chaîne complète composant PDK $\to$ modèle compact $\to$ circuit INTERCONNECT $\to$ contrainte de design en retour sur le layout (annonce du Module~5).
\end{flotfonderie}

\noindent\textbf{Livrable prévu~:} circuit INTERCONNECT utilisant au moins un composant PDK pré-caractérisé + note listant les contraintes de design extraites du PDK utilisé.

\bigskip\hrule\bigskip

\chapter{Optimisation et analyse de tolérance}
\label{chap:m4-tolerance}
\etatunite{PLAN}

\noindent\textbf{Niveau~:} avancé \hfill \textbf{Position~:} Module 4, Chapitre 7 sur 8

\begin{objectifs}
\begin{itemize}[leftmargin=1.4em]
  \item \textbf{[Théorique]} Formaliser la notion de tolérance de fabrication (variation de largeur/épaisseur/indice autour d'une valeur nominale) et son impact statistique sur la performance d'un dispositif.
  \item \textbf{[Pratique]} Utiliser l'outil \emph{sweep/optimization} de Lumerical (ou un script maison via \texttt{lumapi}, cf. \cref{chap:m4-api}) pour explorer l'espace des variations de fabrication du coupleur directionnel.
  \item \textbf{[Pratique]} Produire une analyse de sensibilité (\emph{corner analysis} ou Monte-Carlo simplifié) et en déduire une marge de conception recommandée.
  \item \textbf{[Théorique]} Relier cette analyse à un exercice diagnostic type~: à partir d'un dispositif mesuré hors spécification, remonter à la variation de procédé la plus probable.
\end{itemize}
\end{objectifs}

\begin{prerequis}
\textbf{Théoriques~:} \cref{chap:m4-api,chap:m1-coupleur}.\\
\textbf{Outils/logiciels~:} Lumerical (module \emph{optimization/sweep}) ou script \texttt{lumapi} maison.
\end{prerequis}

\noindent\textbf{Outils principaux~:} Lumerical \emph{sweep/optimization} + \texttt{lumapi} (commercial).

\noindent\textbf{Rappel de mise à niveau prévu~:} non.

\noindent\textbf{Aperçu du contenu~:} chapitre charnière entre simulation \og idéale\fg{} et réalité de fabrication, préparant directement les design rules du Module~5 et le retour mesure$\to$design du Module~6.

\noindent\textbf{Livrable prévu~:} carte de sensibilité performance vs variation de procédé (gap, largeur) + marge de conception recommandée, avec justification chiffrée.

\bigskip\hrule\bigskip

\chapter{Mini-projet Lumerical de synthèse}
\label{chap:m4-projet}
\etatunite{PLAN}

\noindent\textbf{Niveau~:} avancé \hfill \textbf{Position~:} Module 4, Chapitre 8 sur 8

\begin{objectifs}
\begin{itemize}[leftmargin=1.4em]
  \item \textbf{[Pratique]} Concevoir un dispositif complet (MZI thermo-optique ou multiplexeur simple) en combinant FDTD, MODE et INTERCONNECT.
  \item \textbf{[Pratique]} Intégrer une analyse de tolérance (\cref{chap:m4-tolerance}) dans la validation finale du dispositif.
  \item \textbf{[Pratique]} Rédiger une note technique de validation au format attendu d'un ingénieur PIC (hypothèses, résultats, marges, limites).
  \item \textbf{[Théorique]} Faire la synthèse critique des trois méthodes numériques du cours (FDTD/EME/FEM) telles qu'utilisées dans ce projet.
\end{itemize}
\end{objectifs}

\begin{prerequis}
\textbf{Théoriques~:} l'ensemble du Module~4 (\cref{chap:m4-interface,chap:m4-fdtd-lumerical,chap:m4-mode,chap:m4-api,chap:m4-interconnect,chap:m4-pdk,chap:m4-tolerance}).\\
\textbf{Outils/logiciels~:} Lumerical FDTD, MODE, INTERCONNECT, \texttt{lumapi}.
\end{prerequis}

\noindent\textbf{Outils principaux~:} suite Lumerical complète (commercial).

\noindent\textbf{Rappel de mise à niveau prévu~:} non.

\noindent\textbf{Aperçu du contenu~:} projet de synthèse intra-module, distinct du projet de synthèse final du cours (\cref{chap:m7-cahier-des-charges} et suivants) qui, lui, mobilisera en plus le layout/PDK (Module~5) et le flot fonderie (Module~6).

\noindent\textbf{Livrable prévu~:} note technique de validation complète (portfolio) du dispositif conçu, format réutilisable comme modèle pour le projet de synthèse final du Module~7.


% ============================================================
% MODULE 5 — Layout et PDK
% ============================================================
\part{Layout et PDK~: gdsfactory, IPKISS, KLayout, règles de conception fonderie}
\label{part:module5}
% ============================================================
% Module 5 — Layout et PDK
% ÉTAT : PLAN.
% ============================================================

\chapter{Introduction à GDSII et aux outils de layout~: KLayout}
\label{chap:m5-gdsii}
\etatunite{PLAN}

\noindent\textbf{Niveau~:} intermédiaire \hfill \textbf{Position~:} Module 5, Chapitre 1 sur 5

\begin{objectifs}
\begin{itemize}[leftmargin=1.4em]
  \item \textbf{[Théorique]} Décrire le format GDSII (calques/\emph{layers}, cellules, hiérarchie) comme langage commun de description géométrique entre concepteur et fonderie.
  \item \textbf{[Pratique]} Naviguer dans KLayout~: inspection de calques, mesure de distances, navigation hiérarchique d'une cellule.
  \item \textbf{[Pratique]} Ouvrir et annoter un fichier GDS simple (guide droit) généré à partir d'un script minimal.
  \item \textbf{[Théorique]} Relier la notion de calque GDS à la notion de plateforme matériau (\cref{chap:m0-materiaux})~: un calque = une étape de procédé physique.
\end{itemize}
\end{objectifs}

\begin{prerequis}
\textbf{Théoriques~:} \cref{chap:m0-materiaux}.\\
\textbf{Outils/logiciels~:} KLayout (open source, gratuit).
\end{prerequis}

\noindent\textbf{Outils principaux~:} KLayout (open source).

\noindent\textbf{Rappel de mise à niveau prévu~:} non.

\noindent\textbf{Aperçu du contenu~:} première confrontation explicite du cours à un fichier de layout réel, posant le vocabulaire (\emph{layer}, \emph{cell}, \emph{instance}) repris tel quel dans les chapitres gdsfactory/IPKISS suivants.

\noindent\textbf{Livrable prévu~:} fichier GDS annoté (captures KLayout) d'un guide droit simple, avec table de correspondance calque $\to$ matériau/procédé.

\bigskip\hrule\bigskip

\chapter{gdsfactory~: composants paramétriques et génération de layout depuis Python}
\label{chap:m5-gdsfactory}
\etatunite{PLAN}

\noindent\textbf{Niveau~:} intermédiaire \hfill \textbf{Position~:} Module 5, Chapitre 2 sur 5

\begin{objectifs}
\begin{itemize}[leftmargin=1.4em]
  \item \textbf{[Pratique]} Générer par script Python (gdsfactory) les layouts des composants déjà simulés dans le cours~: guide droit (\cref{chap:m1-guide}), coupleur directionnel (\cref{chap:m1-coupleur}), taper (\cref{chap:m2-application}).
  \item \textbf{[Théorique]} Comprendre le principe de PCell (\emph{parametric cell}) et son intérêt pour la traçabilité géométrie $\leftrightarrow$ paramètres de simulation.
  \item \textbf{[Pratique]} Exporter un GDS et le contrôler visuellement dans KLayout (\cref{chap:m5-gdsii}).
  \item \textbf{[Pratique]} Organiser une bibliothèque de composants réutilisable (fonction $\to$ PCell) reprenant les dispositifs déjà validés numériquement dans les Modules 1-4.
\end{itemize}
\end{objectifs}

\begin{prerequis}
\textbf{Théoriques~:} \cref{chap:m5-gdsii,chap:m1-guide,chap:m1-coupleur,chap:m2-application}.\\
\textbf{Outils/logiciels~:} gdsfactory (open source, Python), KLayout.
\end{prerequis}

\noindent\textbf{Outils principaux~:} gdsfactory (open source).

\noindent\textbf{Rappel de mise à niveau prévu~:} non.

\noindent\textbf{Aperçu du contenu~:} bascule de \og géométrie décrite dans le solveur de simulation\fg{} à \og géométrie décrite comme artefact fonderie autonome\fg{}, condition nécessaire au tape-out du Module~6.

\noindent\textbf{Livrable prévu~:} bibliothèque gdsfactory (PCells guide droit, coupleur, taper) + GDS exportés, réutilisés au \cref{chap:m5-drc}.

\bigskip\hrule\bigskip

\chapter{IPKISS~: modélisation paramétrique et co-simulation layout/circuit}
\label{chap:m5-ipkiss}
\etatunite{PLAN}

\noindent\textbf{Niveau~:} avancé \hfill \textbf{Position~:} Module 5, Chapitre 3 sur 5

\begin{objectifs}
\begin{itemize}[leftmargin=1.4em]
  \item \textbf{[Pratique]} Décrire un composant en IPKISS en couplant description géométrique (layout) et modèle de simulation (circuit-level) dans un même objet Python.
  \item \textbf{[Théorique]} Comparer l'approche IPKISS (layout et modèle co-définis) à l'approche gdsfactory + Lumerical vue séparément dans les chapitres précédents (flot découplé).
  \item \textbf{[Pratique]} Reconstruire le MZI simple du \cref{chap:m4-interconnect} en IPKISS et vérifier la cohérence layout/S-paramètres.
  \item \textbf{[Théorique]} Situer IPKISS parmi les outils commerciaux de PDA (\emph{Photonic Design Automation}) et anticiper son rôle dans un flot fonderie complet (Module~6).
\end{itemize}
\end{objectifs}

\begin{prerequis}
\textbf{Théoriques~:} \cref{chap:m5-gdsfactory,chap:m4-interconnect}.\\
\textbf{Outils/logiciels~:} IPKISS (commercial, licence académique à anticiper).
\end{prerequis}

\noindent\textbf{Outils principaux~:} IPKISS (commercial) ; équivalent open source déjà vu~: gdsfactory (\cref{chap:m5-gdsfactory}) pour la partie layout.

\noindent\textbf{Rappel de mise à niveau prévu~:} non.

\noindent\textbf{Aperçu du contenu~:} chapitre de transfert de compétence gdsfactory $\to$ IPKISS, sur le modèle du transfert Meep $\to$ Lumerical déjà pratiqué au \cref{chap:m4-fdtd-lumerical}.

\noindent\textbf{Livrable prévu~:} composant IPKISS du MZI (layout + modèle) + comparaison de la réponse spectrale à celle obtenue en INTERCONNECT (\cref{chap:m4-interconnect}).

\bigskip\hrule\bigskip

\chapter{Règles de conception fonderie~: design rules et DRC photonique}
\label{chap:m5-drc}
\etatunite{PLAN}

\noindent\textbf{Niveau~:} avancé \hfill \textbf{Position~:} Module 5, Chapitre 4 sur 5

\begin{objectifs}
\begin{itemize}[leftmargin=1.4em]
  \item \textbf{[Théorique]} Définir les catégories usuelles de règles de conception (largeur minimale, gap minimal, densité de calque, tolérance de lithographie) et leur origine physique (résolution du procédé de gravure/lithographie).
  \item \textbf{[Pratique]} Exécuter une vérification DRC (\emph{Design Rule Check}) photonique sur les layouts gdsfactory du \cref{chap:m5-gdsfactory} avec KLayout.
  \item \textbf{[Pratique]} Corriger un layout en violation de DRC (ex. gap de coupleur directionnel du \cref{chap:m1-coupleur} trop faible pour la résolution du PDK) et évaluer l'impact sur la performance simulée.
  \item \textbf{[Théorique]} Relier une règle de conception à la tolérance de fabrication déjà quantifiée au \cref{chap:m4-tolerance} — même contrainte physique vue sous deux angles (statistique vs binaire pass/fail).
\end{itemize}
\end{objectifs}

\begin{prerequis}
\textbf{Théoriques~:} \cref{chap:m5-gdsfactory,chap:m4-tolerance}.\\
\textbf{Outils/logiciels~:} KLayout (moteur DRC), fichier de règles DRC (\texttt{.lydrc}) d'un PDK académique/open.
\end{prerequis}

\noindent\textbf{Outils principaux~:} KLayout DRC (open source).

\noindent\textbf{Rappel de mise à niveau prévu~:} non.

\noindent\textbf{Aperçu du contenu~:} premier chapitre du cours où une contrainte fonderie devient un verdict binaire (pass/fail) et non plus seulement une recommandation issue d'une analyse de sensibilité.

\begin{flotfonderie}
Aperçu (à détailler en rédaction complète)~: exemple chiffré (largeur minimale de guide, gap minimal entre guides, tolérance de lithographie typiques d'un PDK MPW) et son impact direct sur le coupleur directionnel du \cref{chap:m1-coupleur}.
\end{flotfonderie}

\noindent\textbf{Livrable prévu~:} rapport DRC (avant/après correction) du coupleur directionnel + note expliquant l'origine physique de chaque règle violée.

\bigskip\hrule\bigskip

\chapter{Vérification LVS photonique et cohérence layout/schématique}
\label{chap:m5-lvs}
\etatunite{PLAN}

\noindent\textbf{Niveau~:} avancé \hfill \textbf{Position~:} Module 5, Chapitre 5 sur 5

\begin{objectifs}
\begin{itemize}[leftmargin=1.4em]
  \item \textbf{[Théorique]} Transposer le concept de LVS (\emph{Layout Versus Schematic}) de la microélectronique à la photonique~: vérifier que le layout GDS correspond bien au circuit INTERCONNECT/IPKISS prévu (connectivité des guides, identité des composants).
  \item \textbf{[Pratique]} Réaliser une vérification LVS photonique (outillage SiEPIC-Tools ou équivalent) sur le MZI du \cref{chap:m5-ipkiss}.
  \item \textbf{[Pratique]} Diagnostiquer une erreur de connectivité typique (guide non jointif, composant mal instancié) introduite volontairement pour l'exercice.
  \item \textbf{[Théorique]} Positionner DRC et LVS comme les deux vérifications obligatoires avant tout dossier de soumission fonderie (transition vers le Module~6).
\end{itemize}
\end{objectifs}

\begin{prerequis}
\textbf{Théoriques~:} \cref{chap:m5-drc,chap:m5-ipkiss}.\\
\textbf{Outils/logiciels~:} SiEPIC-Tools (extension KLayout, open source) ou équivalent disponible.
\end{prerequis}

\noindent\textbf{Outils principaux~:} SiEPIC-Tools / KLayout (open source).

\noindent\textbf{Rappel de mise à niveau prévu~:} non.

\noindent\textbf{Aperçu du contenu~:} clôture du Module~5 sur un flot de vérification complet (DRC + LVS), prérequis direct du dossier de soumission MPW traité au \cref{chap:m6-mpw}.

\noindent\textbf{Livrable prévu~:} rapport LVS du MZI (avec une erreur de connectivité corrigée) + checklist de vérification pré-soumission (DRC + LVS) réutilisée au Module~6.


% ============================================================
% MODULE 6 — Interaction avec la fonderie
% ============================================================
\part{Interaction avec la fonderie~: MPW, DRC/LVS photonique, packaging, itération design-mesure}
\label{part:module6}
% ============================================================
% Module 6 — Interaction avec la fonderie
% ÉTAT : PLAN.
% ============================================================

\chapter{Le flot MPW (Multi-Project Wafer)~: cycle, jalons, coûts, calendrier}
\label{chap:m6-mpw}
\etatunite{PLAN}

\noindent\textbf{Niveau~:} avancé \hfill \textbf{Position~:} Module 6, Chapitre 1 sur 4

\begin{objectifs}
\begin{itemize}[leftmargin=1.4em]
  \item \textbf{[Théorique]} Expliquer le principe du \emph{Multi-Project Wafer}~: mutualisation d'une plaque entre plusieurs projets/clients pour réduire le coût d'un run.
  \item \textbf{[Théorique]} Décrire le cycle typique d'un run MPW (appel à participation, date limite de soumission (\emph{tape-out deadline}), fabrication, retour des puces, délais typiques de plusieurs mois).
  \item \textbf{[Pratique]} Construire un rétroplanning réaliste de projet PIC intégrant les jalons MPW, à partir d'un calendrier de fonderie académique typique (ex. imec/AIM Photonics/CORNERSTONE selon disponibilité).
  \item \textbf{[Théorique]} Identifier les livrables attendus par la fonderie à chaque étape du cycle (pré-soumission, soumission finale, retour de mesure).
\end{itemize}
\end{objectifs}

\begin{prerequis}
\textbf{Théoriques~:} \cref{chap:m5-lvs} (checklist de vérification pré-soumission).\\
\textbf{Outils/logiciels~:} aucun logiciel nouveau ; chapitre organisationnel/méthodologique.
\end{prerequis}

\noindent\textbf{Outils principaux~:} aucun — chapitre de contexte industriel, appuyé sur la documentation publique d'un ou plusieurs \emph{foundry services} académiques.

\noindent\textbf{Rappel de mise à niveau prévu~:} non.

\noindent\textbf{Aperçu du contenu~:} premier chapitre entièrement consacré au vocabulaire et aux processus fonderie (\emph{tape-out}, \emph{shuttle run}, \emph{die}, \emph{reticle}), destiné à un lecteur qui les découvre pour la première fois malgré son expertise physique.

\noindent\textbf{Livrable prévu~:} rétroplanning-type (diagramme simple) d'un projet PIC académique, des premières simulations à la réception des puces.

\bigskip\hrule\bigskip

\chapter{Dossier de soumission fonderie~: documentation, PDK, checklist de tape-out}
\label{chap:m6-dossier}
\etatunite{PLAN}

\noindent\textbf{Niveau~:} avancé \hfill \textbf{Position~:} Module 6, Chapitre 2 sur 4

\begin{objectifs}
\begin{itemize}[leftmargin=1.4em]
  \item \textbf{[Pratique]} Constituer un dossier de soumission complet~: GDS final vérifié (DRC/LVS, \cref{chap:m5-drc,chap:m5-lvs}), fiche de calques utilisés, description des composants et de leur origine (PDK vs custom).
  \item \textbf{[Théorique]} Distinguer les responsabilités respectives du concepteur et de la fonderie dans la validation finale d'un design (la fonderie ne revalide pas la physique, seulement la conformité géométrique/procédé).
  \item \textbf{[Pratique]} Remplir une checklist de tape-out type à partir des éléments déjà produits dans le cours (guide, coupleur, taper, MZI).
  \item \textbf{[Théorique]} Anticiper les causes usuelles de rejet d'un dossier de soumission (violations DRC non corrigées, calques manquants, métadonnées absentes).
\end{itemize}
\end{objectifs}

\begin{prerequis}
\textbf{Théoriques~:} \cref{chap:m6-mpw,chap:m5-lvs}.\\
\textbf{Outils/logiciels~:} KLayout (assemblage final du GDS), gabarits de documentation (à fournir en rédaction complète).
\end{prerequis}

\noindent\textbf{Outils principaux~:} KLayout (open source) pour l'assemblage final ; pas d'outil commercial nouveau.

\noindent\textbf{Rappel de mise à niveau prévu~:} non.

\noindent\textbf{Aperçu du contenu~:} chapitre pratique de synthèse documentaire, transformant les livrables épars des Modules 1-5 en un dossier de soumission unique et cohérent.

\noindent\textbf{Livrable prévu~:} dossier de soumission complet (GDS + documentation) pour le MZI du \cref{chap:m5-ipkiss}, prêt (au format) pour un run MPW académique.

\bigskip\hrule\bigskip

\chapter{Packaging et interfaces optiques/électriques}
\label{chap:m6-packaging}
\etatunite{PLAN}

\noindent\textbf{Niveau~:} avancé \hfill \textbf{Position~:} Module 6, Chapitre 3 sur 4

\begin{objectifs}
\begin{itemize}[leftmargin=1.4em]
  \item \textbf{[Théorique]} Comparer les deux stratégies principales d'interface optique fibre-puce~: couplage par la tranche (\emph{edge coupling}) et couplage par réseau de diffraction (\emph{grating coupler}), en s'appuyant sur les notions de mode déjà acquises (\cref{chap:m0-optique}).
  \item \textbf{[Pratique]} Simuler (rapidement, en réutilisant les outils du Module~1 ou 4) l'efficacité de couplage d'un \emph{grating coupler} simple et identifier ses paramètres critiques.
  \item \textbf{[Théorique]} Décrire les interfaces électriques usuelles pour les composants actifs (\emph{wire bonding}, \emph{flip-chip}) et leur impact sur l'encombrement et le placement des \emph{pads}.
  \item \textbf{[Théorique]} Relier le choix d'interface de packaging à des contraintes de layout déjà rencontrées (Module~5)~: zones d'exclusion, marges de découpe (\emph{dicing}).
\end{itemize}
\end{objectifs}

\begin{prerequis}
\textbf{Théoriques~:} \cref{chap:m0-optique,chap:m6-dossier}.\\
\textbf{Outils/logiciels~:} réutilisation de Meep ou Lumerical FDTD (\cref{chap:m1-meep,chap:m4-fdtd-lumerical}) pour la simulation du grating coupler.
\end{prerequis}

\noindent\textbf{Outils principaux~:} Meep ou Lumerical FDTD (réutilisation d'un outil déjà maîtrisé, pas de nouvel outil dans ce chapitre).

\noindent\textbf{Rappel de mise à niveau prévu~:} oui, court rappel sur le principe de diffraction d'un réseau, par analogie avec les réseaux de diffraction déjà rencontrés en optique classique.

\noindent\textbf{Aperçu du contenu~:} chapitre qui referme la boucle simulation $\to$ layout $\to$ contrainte physique de packaging, dernier maillon avant la mesure.

\noindent\textbf{Livrable prévu~:} courbe d'efficacité de couplage d'un grating coupler simple vs longueur d'onde + note de comparaison edge coupling / grating coupler pour le cas d'usage du cours.

\bigskip\hrule\bigskip

\chapter{Boucle design-mesure-itération}
\label{chap:m6-boucle}
\etatunite{PLAN}

\noindent\textbf{Niveau~:} avancé \hfill \textbf{Position~:} Module 6, Chapitre 4 sur 4

\begin{objectifs}
\begin{itemize}[leftmargin=1.4em]
  \item \textbf{[Théorique]} Décrire le principe général de caractérisation post-fabrication d'un PIC (banc de mesure fibre-à-fibre, sources large bande, détecteurs) sans viser l'exhaustivité métrologique.
  \item \textbf{[Pratique]} Comparer un spectre de mesure simulé (généré volontairement avec un écart réaliste) au spectre simulé initial (Module~4) pour illustrer l'écart design-mesure typique.
  \item \textbf{[Pratique]} Exercice diagnostic de synthèse~: à partir d'un écart mesure/simulation donné, remonter à la cause la plus probable en mobilisant l'ensemble du cours (maillage, tolérance de fabrication, condition aux limites, erreur de packaging).
  \item \textbf{[Théorique]} Formaliser la boucle d'itération design $\to$ fabrication $\to$ mesure $\to$ correction de modèle, et son coût temporel (plusieurs mois par itération en MPW académique).
\end{itemize}
\end{objectifs}

\begin{prerequis}
\textbf{Théoriques~:} l'ensemble du Module~6 et, en pratique, l'essentiel du cours (\cref{chap:m4-tolerance,chap:m6-packaging}).\\
\textbf{Outils/logiciels~:} réutilisation des outils Python/Lumerical déjà en place ; aucun nouvel outil.
\end{prerequis}

\noindent\textbf{Outils principaux~:} aucun nouvel outil — chapitre de synthèse méthodologique.

\noindent\textbf{Rappel de mise à niveau prévu~:} non.

\noindent\textbf{Aperçu du contenu~:} chapitre culminant du Module~6, formulé comme un exercice de type support technique client (FAE) mobilisant transversalement les Modules 1 à 6, en préparation directe du projet de synthèse final (Module~7).

\begin{diagnostic}[Exercice diagnostic de synthèse (aperçu)]
À détailler en rédaction complète~: un spectre de mesure d'un MZI présente un décalage de longueur d'onde de résonance et une extinction dégradée par rapport à la simulation device-level du \cref{chap:m4-projet}. Identifier, à partir des seuls éléments fournis (écart chiffré, tolérances connues du PDK), la ou les causes les plus probables parmi~: variation de largeur de guide hors tolérance, erreur de calibration du grating coupler, dérive thermique non modélisée.
\end{diagnostic}

\noindent\textbf{Livrable prévu~:} note de diagnostic complète (cause identifiée, justification, recommandation de correction de design), clôturant le Module~6.


% ============================================================
% MODULE 7 — Projet de synthèse
% ============================================================
\part{Projet de synthèse~: dispositif complet, simulation multi-outils, note d'application finale}
\label{part:module7}
% ============================================================
% Module 7 — Projet de synthèse
% ÉTAT : PLAN.
% ============================================================

\chapter{Cahier des charges d'un dispositif complet}
\label{chap:m7-cahier-des-charges}
\etatunite{PLAN}

\noindent\textbf{Niveau~:} avancé \hfill \textbf{Position~:} Module 7, Chapitre 1 sur 3

\begin{objectifs}
\begin{itemize}[leftmargin=1.4em]
  \item \textbf{[Pratique]} Rédiger un cahier des charges technique complet pour un dispositif de synthèse (choix proposé~: MZI thermo-optique reconfigurable \emph{ou} (dé)multiplexeur WDM 2 canaux — à trancher en tête de chapitre selon l'équilibre théorie/pratique visé).
  \item \textbf{[Théorique]} Décomposer le dispositif choisi en sous-composants déjà maîtrisés individuellement dans le cours (guide, coupleur/MZI, taper, éventuellement élément thermo-optique du \cref{chap:m3-multiphysique}).
  \item \textbf{[Pratique]} Définir des spécifications quantitatives mesurables (bande passante, extinction, pertes d'insertion, budget de longueur imposé par le PDK) servant de critère d'acceptation final.
  \item \textbf{[Théorique]} Identifier, pour chaque spécification, quelle méthode du cours (FDTD/EME/FEM/circuit-level) est la mieux adaptée pour la vérifier.
\end{itemize}
\end{objectifs}

\begin{prerequis}
\textbf{Théoriques~:} l'ensemble des Modules 0 à 6 (chapitre de synthèse).\\
\textbf{Outils/logiciels~:} aucun nouveau ; chapitre de cadrage méthodologique.
\end{prerequis}

\noindent\textbf{Outils principaux~:} aucun — production d'un document de spécification.

\noindent\textbf{Rappel de mise à niveau prévu~:} non.

\noindent\textbf{Aperçu du contenu~:} chapitre d'entrée du projet final, conçu comme exercice de \og traduction\fg{} d'un besoin fonctionnel en plan de simulation multi-outils, compétence transversale visée par l'ensemble du cours.

\noindent\textbf{Livrable prévu~:} cahier des charges rédigé (spécifications quantitatives + plan de vérification par méthode), validé avant simulation au \cref{chap:m7-simulation}.

\bigskip\hrule\bigskip

\chapter{Simulation multi-outils intégrée et validation croisée}
\label{chap:m7-simulation}
\etatunite{PLAN}

\noindent\textbf{Niveau~:} avancé \hfill \textbf{Position~:} Module 7, Chapitre 2 sur 3

\begin{objectifs}
\begin{itemize}[leftmargin=1.4em]
  \item \textbf{[Pratique]} Simuler chaque sous-composant du dispositif avec l'outil le plus approprié (FDTD pour la validation locale, EME/MODE pour les tapers, FEM pour un cas géométrique complexe si pertinent, circuit-level pour l'assemblage final).
  \item \textbf{[Pratique]} Assembler les résultats en un circuit INTERCONNECT/IPKISS complet (\cref{chap:m4-interconnect,chap:m5-ipkiss}) et vérifier la conformité aux spécifications du \cref{chap:m7-cahier-des-charges}.
  \item \textbf{[Pratique]} Intégrer une analyse de tolérance (\cref{chap:m4-tolerance}) et une vérification DRC/LVS (\cref{chap:m5-drc,chap:m5-lvs}) dans la validation finale.
  \item \textbf{[Théorique]} Documenter les écarts entre méthodes lorsqu'ils apparaissent, en s'appuyant sur les réflexes diagnostiques développés tout au long du cours.
\end{itemize}
\end{objectifs}

\begin{prerequis}
\textbf{Théoriques~:} \cref{chap:m7-cahier-des-charges} et l'ensemble des outils du cours (Modules 1 à 5).\\
\textbf{Outils/logiciels~:} Meep, Lumerical (FDTD/MODE/INTERCONNECT), gdsfactory/IPKISS, KLayout — sélection selon les choix du \cref{chap:m7-cahier-des-charges}.
\end{prerequis}

\noindent\textbf{Outils principaux~:} l'ensemble de la boîte à outils du cours, mobilisée simultanément pour la première fois.

\noindent\textbf{Rappel de mise à niveau prévu~:} non.

\noindent\textbf{Aperçu du contenu~:} cœur pratique du projet de synthèse~; chapitre le plus long du Module~7, organisé sous-composant par sous-composant puis assemblage final, à l'image de la structure \og mise en place $\to$ paramètres $\to$ exécution $\to$ extraction $\to$ interprétation\fg{} employée dans tout le cours.

\noindent\textbf{Livrable prévu~:} dossier de simulation complet du dispositif (tous sous-composants + assemblage circuit-level + vérifications DRC/LVS/tolérance), prêt pour rédaction de la note finale.

\bigskip\hrule\bigskip

\chapter{Note d'application finale et portfolio technique}
\label{chap:m7-note-finale}
\etatunite{PLAN}

\noindent\textbf{Niveau~:} avancé \hfill \textbf{Position~:} Module 7, Chapitre 3 sur 3

\begin{objectifs}
\begin{itemize}[leftmargin=1.4em]
  \item \textbf{[Pratique]} Rédiger une note d'application technique complète (format proche d'une \emph{application note} industrielle)~: contexte, méthode, résultats, marges, limites, recommandations.
  \item \textbf{[Pratique]} Assembler l'ensemble des livrables du cours (Modules 0 à 7) en un portfolio technique cohérent, démontrant la maîtrise du flot simulation $\to$ layout $\to$ fonderie.
  \item \textbf{[Théorique]} Réaliser une synthèse critique personnelle~: pour chaque méthode (FDTD/EME/FEM) et chaque outil (Meep/Lumerical/IPKISS/gdsfactory/KLayout), formuler en une phrase son cas d'usage optimal.
  \item \textbf{[Théorique]} Identifier, en conclusion, les sujets hors périmètre du cours qui constitueraient une suite logique (ex. photonique non linéaire intégrée, calcul quantique photonique, co-intégration électronique-photonique avancée).
\end{itemize}
\end{objectifs}

\begin{prerequis}
\textbf{Théoriques~:} \cref{chap:m7-simulation}.\\
\textbf{Outils/logiciels~:} aucun nouveau ; chapitre de rédaction et de synthèse.
\end{prerequis}

\noindent\textbf{Outils principaux~:} aucun — production documentaire finale.

\noindent\textbf{Rappel de mise à niveau prévu~:} non.

\noindent\textbf{Aperçu du contenu~:} chapitre de clôture du cours, sans contenu théorique nouveau, entièrement dédié à la consolidation et à la valorisation (portfolio) des livrables produits depuis le Module~0.

\noindent\textbf{Livrable prévu~:} note d'application finale complète + portfolio technique consolidé (tous livrables du cours), livrable ultime du cours dans son ensemble.


\end{document}
